\begin{abstract}

Questo lavoro presenta un tool per la ricerca automatica di virtual patients in reti biochimiche, basato sulla formulazione del problema come ottimizzazione vincolata black-box.
A partire da modelli qualitativi ottenuti dal database Reactome, il tool effettua una conversione automatica in modelli quantitativi SBML, aggiungendo leggi cinetiche, parametri e condizioni iniziali.
La stima dei parametri è formalizzata come un problema di minimizzazione:

\begin{equation}
\hat{\theta} = \argmin_{\theta} \mathcal{L}(\theta)
\end{equation}

Dove $\mathcal{L}$ combina più termini di perdita per garantire la stabilità dinamica, la normalizzazione delle concentrazioni e la plausibilità biologica dell'ordinamento delle specie.
L'ottimizzazione è vincolata da relazioni sulle costanti cinetiche e da limiti sulle concentrazioni.
Il framework integra diversi algoritmi evolutivi e numerici (Nomad, Nevergrad, e un algoritmo genetico OpenAI) per esplorare lo spazio dei parametri e identificare configurazioni compatibili con i vincoli biologici.
I risultati mostrano che il sistema è in grado di discriminare ordinamenti non plausibili e individuare set di parametri coerenti con le osservazioni biologiche, fornendo una base automatizzata per la simulazione di modelli biologici.

Viene mostrato l'utilizzo del tool su 5 modelli presi da Reactome.

\end{abstract}